%------------------------------------------------------------------------------%

\usepackage{apresentacao}
\usepackage{mathconfig}
\usepackage{tabto}

\makeatletter
\graphicspath  {{./figs/}}
\def\input@path{{./figs/}}
\makeatother

\newcounter{exemplo}
\setcounter{exemplo}{0}

%------------------------------------------------------------------------------%
\title {Séries Numéricas com Termos Gerais}
\author{\large Luis Alberto D'Afonseca}
\date  {Integrais e Séries}

%------------------------------------------------------------------------------%
\begin{document}

%------------------------------------------------------------------------------%
\begin{frame}
    \titlepage
\end{frame}

%------------------------------------------------------------------------------%
\section{Séries Alternadas}

%------------------------------------------------------------------------------%
\begin{frame}{Séries Alternadas}
  Os \emph{termos trocam de sinal}
  \pause
  alternando entre um positivo e um negativo
  \[ \sum_{n=1}^\infty (-1)^{n+1} a_n \]
  com $a_n >0$ para todo $n$

  \pause
  \vspace{\baselineskip}
  Exemplo, \emph{Série Harmônica Alternada}
  \[
    \sum_{n=1}^\infty \frac{(-1)^{n+1}}{n}
    \;=\; 1 - \frac{1}{2} + \frac{1}{3} - \frac{1}{4} + \cdots
  \]
\end{frame}

%------------------------------------------------------------------------------%
\begin{frame}{Teste da Série Alternada -- Teste de Leibniz}
  \onslide<+->{
  Seja $\,a_n\,$ tal que, para todo $ n \in \N$,}

  \vspace{0.5\baselineskip}
  \begin{enumerate}[<+->]
      \item $a_n>0$           \tabto{30mm} $a_n$ são os módulos dos termos da série  \\[2mm]
      \item $a_n\geq a_{n+1}$ \tabto{30mm} os termos são decrescentes em módulo      \\[2mm]
      \item $a_n\to 0$        \tabto{30mm} os termos tendem a zero
  \end{enumerate}

  \onslide<+->{
  \vspace{\baselineskip}
  Então, a série alternada
  \(\; \sum_{n=1}^\infty (-1)^{n+1}a_n \;\)
  converge}
\end{frame}

%------------------------------------------------------------------------------%
\addtocounter{exemplo}{1}
\begin{frame}{Exemplo \theexemplo}
  \onslide<+->{Vamos analisar a série
  \[
    \sum_{n=1}^{\infty} (-1)^{n+1} \frac{1}{\,2^{n-1}\,}
  \]}
  \onslide<2->{Percebemos que}
  \begin{itemize}[<+->]
    \item a série é alternada
    \item o módulo do termo geral é \(\quad a_n = \frac{1}{\,2^{n-1}\,} \)
    \item $a_n\geq a_{n+1}$
    \item $a_n\to 0$
  \end{itemize}
  \onslide<+->{Portanto a série converge pelo Teste da Série Alternada}
\end{frame}

%------------------------------------------------------------------------------%
\begin{frame}{Exemplo \theexemplo}
  Por outro lado
  \[
    \sum_{n=1}^{\infty} (-1)^{n+1} \frac{1}{\,2^{n-1}\,}
    \pause
    \;=\; \sum_{n=1}^{\infty} \frac{(-1)^2(-1)^{n-1}}{\,2^{n-1}\,}
    \pause
    \;=\; \sum_{n=1}^{\infty} \left( -\frac{1}{\,2\,} \right)^{n-1}
  \]
  \pause
  Série geométrica com $\;a=1\;$ e \(\;r=-\frac{1}{2}\;\) \\[5mm]

  \pause
  \vspace{0.5\baselineskip}
  Como $|r|<1$ ela converge \\[5mm]

  \pause
  sua soma é
  \(\quad
    S \;=\; \frac{a}{1-r}
      \pause
      \;=\; \frac{1}{1+\sfrac{1}{2}}
      \;=\; \frac{2}{3}
  \)

\end{frame}

%------------------------------------------------------------------------------%
%\framedgraphic{Usando Planilhas Eletrônicas}{./figs/exemplo_planilha}{}
%\framedgraphic{Usando Planilhas Eletrônicas}{./figs/exemplo_grafico}{}

%------------------------------------------------------------------------------%
\section{Estimativa do Erro para Séries Alternadas}

%------------------------------------------------------------------------------%
\begin{frame}{Estimativa do Erro}
  Dada uma série que satisfaz as condições do teste da série alternada
  \pause
  \[ S_n = a_1 - a_2 + a_3 - a_4 + \cdots + (-1)^{n+1}a_n \]
  converge para a soma da série $S$
  \pause
  \begin{itemize}[<+->]
    \item \( \abs{R_n} < a_{n+1} \)
    \item $S$ está entre $S_n$ e $S_{n+1}$
    \item \( R_n = S - S_n \;\) tem o mesmo sinal do primeiro termo não utilizado
  \end{itemize}
\end{frame}

%------------------------------------------------------------------------------%
\section{Convergência Absoluta}

%------------------------------------------------------------------------------%
\begin{frame}{Convergência Absoluta}
  Uma série
  \(\quad \sum_{n=1}^{\infty}\; a_n\) \\[5mm]
  \pause
  \emph{converge absolutamente}, ou é \emph{absolutamente convergente}, \\[5mm]
  \pause
  se
  \(\quad\sum_{n=1}^{\infty} \;\abs{a_n} \quad\) converge
\end{frame}

%------------------------------------------------------------------------------%
\begin{frame}{Convergência Condicional}
  Se  \tabto{12mm} \(\sum_{n=1}^{\infty} a_n\)          \tabto{32mm} \emph{converge} \\[9mm]
  \pause
  mas \tabto{12mm} \(\sum_{n=1}^{\infty} \;\abs{a_n} \) \tabto{32mm} \emph{diverge}  \\[9mm]
  \pause
  a série é \emph{condicionalmente convergente} \\[9mm]
  \pause
  Exemplo: Série Harmônica Alternada
\end{frame}

%------------------------------------------------------------------------------%
\begin{frame}{Teste da Convergência Absoluta}
  Uma série absolutamente convergente é convergente.\\[3mm]

  \pause
  \vspace{\baselineskip}
  Séries com termos não negativos e convergentes
  são absolutamente convergentes
\end{frame}

%------------------------------------------------------------------------------%
\begin{frame}{Teorema do Rearranjo}
Se \(\quad\sum_{n=1}^{\infty} a_n \quad \) converge absolutamente \\[5mm]

\pause
dado qualquer \emph{rearranjo} $(b_k)$ da sequência $(a_n)$ \\[5mm]

\pause
a série $\sum b_k$ converge e
\[
  \sum_{k=1}^{\infty} b_k = \sum_{n=1}^{\infty} a_n
\]
\end{frame}

%%------------------------------------------------------------------------------%
%%------------------------------------------------------------------------------%
%\section{Resumo dos Testes}
%%------------------------------------------------------------------------------%
%%------------------------------------------------------------------------------%
%
%%------------------------------------------------------------------------------%
%\begin{frame}{Testes Básicos}
%	\begin{itemize}[<+->]
%		\item \emph{Teste da divergência ou do $n$-ésimo termo}\\
%          a menos que $a_n\to 0$, a série diverge \\[5mm]
%		\item \emph{Séries geométricas} \\[3mm]
%          \(\sum_{n=1}^{\infty} ar^{n-1}\) converge se \(\abs{r}<1\); caso contrário, diverge \\[5mm]
%		\item \emph{Série $p$} \\[3mm]
%          \(\sum_{n=1}^{\infty} \frac{1}{n^p}\)
%          converge se $p>1$; caso contrário, diverge\\[5mm]
%    \item Série \emph{Telescópica}
%	\end{itemize}
%\end{frame}
%
%%------------------------------------------------------------------------------%
%\begin{frame}{Séries de Termos Não-Negativos}
%\begin{itemize}[<+->]
%  \item Teste da \emph{integral} \\[5mm]
%  \item Teste da \emph{razão} ou o teste da \emph{raiz} \\[5mm]
%  \item Teste da \emph{comparação}
%\end{itemize}
%\end{frame}
%
%%------------------------------------------------------------------------------%
%\begin{frame}{Séries Gerais -- Alguns Termos Negativos}
%	\begin{itemize}[<+->]
%    \item A série é \emph{alternada}? \\[3mm]
%          Verifique as condições do teste da série alternada \\[5mm]
%		\item \(\sum_{n=1}^{\infty} \abs{a_n}\) converge? \\[3mm]
%		      Se convergir, \(\sum_{n=1}^{\infty} a_n\) também converge
%	\end{itemize}
%\end{frame}

%------------------------------------------------------------------------------%
\section{Exemplos}

%------------------------------------------------------------------------------%
\addtocounter{exemplo}{1}
\begin{frame}{Exemplo \theexemplo}
  Use o Teste de Leibniz para mostrar que a série a seguir converge
  \[
    \sum_{n=1}^\infty(-1)^{n}\ln\left(1+\frac{1}{n}\right)
  \]
\end{frame}

%------------------------------------------------------------------------------%
\begin{frame}{Exemplo \theexemplo\ -- Condição 1: \(a_n>0\)}
  \pause
  Como
  \(
    1 + \frac{1}{n} > 1
  \)
  para todo $n\geq 1$

  \pause
  \vspace{\baselineskip}
  Podemos escolher
  \[
    a_n = \ln\left(1+\frac{1}{n}\right)
  \]
  que garante $a_n>0$
\end{frame}

%------------------------------------------------------------------------------%
\begin{frame}{Exemplo \theexemplo\ -- Condição 2: \(a_n\geq a_{n+1}\)}
  \pause
  Considerando a função
  \[
    f(x)
    = \ln\left(1+\frac{1}{x}\right)
    \pause
    = \ln\left(1+x^{-1}\right)
  \]
  \pause
  temos
  \[
    f'(x)
    = \frac{1}{1+x^{-1}} \, \left(-x^{-2}\right)
    \pause
    = \frac{-1}{\left(1+x^{-1}\right)x^2}
    \pause
    = \frac{-1}{x^2+x}
    \pause
    < 0
    \qquad \forall x\geq 1
  \]

  \pause
  \vspace{0.5\baselineskip}
  Portanto $a_n$ é decrescente \(a_n>a_{n+1}\)
\end{frame}

%------------------------------------------------------------------------------%
\begin{frame}{Exemplo \theexemplo\ -- Condição 3: \(a_n\to 0\)}
  \pause
  Observando que
  \[
    \lim_{n\to\infty} 1+\frac{1}{n} = 1
  \]
  \pause
  e a função $\ln(x)$ é contínua em $1$
  \pause
  \[
    \lim a_n
    = \lim \ln \left(1+\frac{1}{n}\right)
    \pause
    = \ln\lim \left(1+\frac{1}{n}\right)
    \pause
    = \ln 1
    \pause
    = 0
  \]

  \pause
  O Teste de Leibniz garante que a série converge
\end{frame}
%------------------------------------------------------------------------------%

%------------------------------------------------------------------------------%
\addtocounter{exemplo}{1}
\begin{frame}{Exemplo \theexemplo}
  Use o Teste de Leibniz para mostrar que a série a seguir converge
  \[
    \sum_{n=2}^\infty(-1)^{n+1}\frac{4}{(\ln n)^2}
  \]
\end{frame}

%------------------------------------------------------------------------------%
\begin{frame}{Exemplo \theexemplo\ -- Condição 1: \(a_n>0\)}
  \pause
  Vemos que
  \(
    a_n = \frac{4}{(\ln n)^2} > 0
  \)
  para todo $n\geq 2$
\end{frame}

%------------------------------------------------------------------------------%
\begin{frame}{Exemplo \theexemplo\ -- Condição 2: \(a_n\geq a_{n+1}\)}
  \pause
  Como  \(\ln n\) é crescente\pause,
  \[
    a_n = \frac{4}{(\ln n)^2}
  \]
  é decrescente
\end{frame}

%------------------------------------------------------------------------------%
\begin{frame}{Exemplo \theexemplo\ -- Condição 3: \(a_n\to 0\)}
  \pause
  $a_n\to 0$, pois
  \[
    \left(\ln n\right)^2\to \infty
  \]

  \pause
  \vspace{\baselineskip}
  O teste de Leibniz garante que a série converge
\end{frame}

%------------------------------------------------------------------------------%
\addtocounter{exemplo}{1}
\begin{frame}{Exemplo \theexemplo}
  Use o teste de Leibniz para mostrar que a série a seguir converge
  \[
    \sum_{n=1}^\infty(-1)^{n+1}\frac{10^n}{(n+1)!}
  \]
\end{frame}

%------------------------------------------------------------------------------%
\begin{frame}{Exemplo \theexemplo\ -- Condição 1: \(a_n>0\)}
  \pause
  \(
    a_n = \frac{10^n}{(n+1)!}
  \)
  é positivo para \(n\geq 1\)
\end{frame}

%------------------------------------------------------------------------------%
\begin{frame}{Exemplo \theexemplo\ -- Condição 2: \(a_n\geq a_{n+1}\)}
  \begin{columns}
    \column{0.5\linewidth}
    \begin{align*}
      \onslide<+->{n   & \geq 8  \\[2mm]}
      \onslide<+->{n+2 & \geq 10 \\[2mm]}
      \onslide<+->{1   & \geq \frac{10}{n+2}}
    \end{align*}
    \column{0.5\linewidth}
    \begin{align*}
      \onslide<+->{a_{n+1}
                   & =\frac{10^{n+1}}{(n+2)!}             \\}
      \onslide<+->{& =\frac{10}{(n+2)}\frac{10^n}{(n+1)!} \\}
      \onslide<+->{& \leq \frac{10^n}{(n+1)!}             \\}
      \onslide<+->{& =a_n}
    \end{align*}

    \onslide<+->{
      Portanto, $a_n$ é decrescente para \(n\geq8\)}
  \end{columns}
\end{frame}

%------------------------------------------------------------------------------%
\begin{frame}{Exemplo \theexemplo\ -- Condição 3: \(a_n\to 0\)}
  \onslide<+->{}
  \onslide<+->{Para \(n>10\)}
  \begin{align*}
    \onslide<+->{(n+1)!
      & = 1 \times 2 \times \cdots
            \times 9 \times 10 \times 11 \times \cdots
            \times n \times (n+1) \\}
    \onslide<+->{& \geq 10! \; 10^{n-10}\,(n+1)}
  \end{align*}
  \onslide<+->{Dessa forma}
  \[
  \onslide<5->{\frac{10^n}{(n+1)!}
    \leq \frac{10^n}{10!\;10^{n-10}\,(n+1)}}
  \onslide<+->{=\frac{10^{10}}{10!}\,\frac{1}{(n+1)}}
  \onslide<+->{\to 0}
  \]

  \onslide<+->{Como
    \(
      0 <
      \frac{10^n}{(n+1)!}
      \leq \frac{10^{10}}{10!}\,\frac{1}{(n+1)}
    \)\\[5mm]}

  \onslide<+->{o Teorema do Confronto garante que $a_n\to 0$}
\end{frame}

%------------------------------------------------------------------------------%
\begin{frame}{Exemplo \theexemplo}
  O teste de Leibniz garante que a série converge
\end{frame}

%------------------------------------------------------------------------------%
\addtocounter{exemplo}{1}
\begin{frame}{Exemplo \theexemplo}
  Estime o erro cometido quando aproximamos o valor da soma da série
  \[
    \sum_{n=1}^\infty \frac{(-1)^{n+1}}{n}
  \]
  pela soma dos seus quatro primeiros termos
\end{frame}

%------------------------------------------------------------------------------%
\begin{frame}{Exemplo \theexemplo}
  A série converge pelo teste da série alternada, portanto
  \[
    \abs{R_n} < a_{n+1}
  \]

  \pause
  Como somamos quatro quatro termos (\(S_4\)), temos que
  \[
    \abs{R_4}
    < a_5
    \pause
    = \frac{1}{5}
  \]
\end{frame}

%------------------------------------------------------------------------------%
\addtocounter{exemplo}{1}
\begin{frame}{Exemplo \theexemplo}
  Determine quantos termos devem ser somados para aproximar a soma da série
  \[
    \sum_{n=1}^\infty \frac{(-1)^{n}}{n^2+3}
  \]
  com erro menor do que 0,0001
\end{frame}

%------------------------------------------------------------------------------%
\begin{frame}{Exemplo \theexemplo}
  \onslide<+->{A série converge pelo teste da série alternada, portanto
  \[
    \abs{R_n} < a_{n+1}
  \]}

  \onslide<+->{
  Para garantir que o erro é menor do que 0,001 devemos impor que}
\begin{columns}
\column{0.5\linewidth}
  \begin{align*}
    \onslide<2->{a_{n+1}           & < 10^{-4} \\[2mm]}
    \onslide<+->{\frac{1}{n^2+3}   & < 10^{-4} \\[2mm]}
    \onslide<+->{\frac{1}{10^{-4}} & < n^2+3}
  \end{align*}
\column{0.5\linewidth}
  \begin{align*}
    \onslide<+->{n^2 & > 10^4 -3 \\[2mm]}
    \onslide<+->{n^2 & > 10^4    \\[2mm]}
    \onslide<+->{n   & > 10^2 }
  \end{align*}
\end{columns}
\end{frame}

%------------------------------------------------------------------------------%
\begin{frame}{\contentsname}
  \tableofcontents
\end{frame}

%------------------------------------------------------------------------------%
\end{document}
%------------------------------------------------------------------------------%
